
\section{Presentación del sistema}

Netflix es una plataforma de entretenimiento por Internet que ofrece series, películas, videojuegos y programación en directo. Sus usuarios pueden acceder a estos contenidos desde diferentes dispositivos mediante su servicio de streaming.

La empresa tiene presencia internacional y organiza sus resultados de streaming en cuatro regiones: Estados Unidos y Canadá, Europa, Medio Oriente y África, Latinoamérica y Asia-Pacífico.

Según su informe anual de 2025, Netflix registró ingresos por 45.183 millones de dólares, frente a los 39.001 millones obtenidos en 2024. Esto representó un crecimiento del 16\% en un año. La empresa relaciona este aumento principalmente con el crecimiento de las membresías, los cambios en los precios y el incremento de los ingresos por publicidad.

Estas cifras permiten comprender la escala económica de sus operaciones y su alcance internacional. Para ofrecer sus contenidos a usuarios de diferentes regiones y dispositivos, Netflix necesita una infraestructura distribuida que permita entregar el contenido y mantener la comunicación entre los componentes de su plataforma.

Por estas características, Netflix es un caso interesante para estudiar cómo se organiza un sistema distribuido real.

\parencite{netflix2025annual}
